\clearpage\phantomsection\label{indices-debut}\pdfbookmark[0]{Indications des défis}{bm-indices}
\editionPage{indices-1}{Indications progressives}{Une aide, puis une nouvelle recherche}{Ne consulter que l’aide nécessaire.}
\begin{blocEdition}{indices-1}
\editionEtape{D01 · Décimales finies et périodes}
\textbf{Indice 1.} Une écriture décimale finie s’écrit $N/10^k$ ; utiliser $10^k=2^k5^k$.\par\smallskip
\textbf{Indice 2.} Dans une division par $q$, les restes possibles sont $0,1,\ldots,q-1$. Que se passe-t-il si un reste revient ?\par
\editionEtape{D02 · Deux identités à construire}
\textbf{Indice 1.} Traiter $a+b+c$ comme une somme de trois termes et distribuer chaque terme.\par\smallskip
\textbf{Indice 2.} Dans le cube, regrouper les trois occurrences de $a^2b$ et les trois de $ab^2$.\par
\editionEtape{D03 · Moyennes et rectangle optimal}
\textbf{Indice 1.} Développer le carré puis isoler $2\sqrt{ab}$.\par\smallskip
\textbf{Indice 2.} Pour la dernière question, $\sqrt{ab}$ est fixe ; l’inégalité donne une borne de $a+b$.\par
\editionEtape{D04 · Barycentres de deux et trois points}
\textbf{Indice 1.} Écrire $\vect{GA}=A-G$ en coordonnées, puis regrouper les termes contenant $G$.\par\smallskip
\textbf{Indice 2.} Des poids positifs donnent un point entre deux points ; un poids négatif peut produire un point extérieur au segment.\par

\end{blocEdition}
\editionPage{indices-2}{Indications progressives}{Une aide, puis une nouvelle recherche}{Ne consulter que l’aide nécessaire.}
\begin{blocEdition}{indices-2}
\editionEtape{D05 · Pourquoi les médianes se coupent}
\textbf{Indice 1.} Remplacer $I$ par $(B+C)/2$ et comparer les expressions de $G-A$ et de $I-A$.\par\smallskip
\textbf{Indice 2.} Le calcul ne privilégie aucun sommet : reprendre le même argument depuis $B$ et $C$.\par
\editionEtape{D06 · Concours des hauteurs}
\textbf{Indice 1.} Additionner les coordonnées de $B'$ et de $C'$, puis diviser par 2.\par\smallskip
\textbf{Indice 2.} La hauteur passant par $A$ devient la médiatrice d’un côté du grand triangle.\par
\editionEtape{D07 · Une aire à partir de deux côtés et d’un angle}
\textbf{Indice 1.} Prendre $CB$ comme base et abaisser la hauteur depuis $A$.\par\smallskip
\textbf{Indice 2.} Les angles supplémentaires ont ici le même sinus ; cela ne fixe pas la forme complète du triangle.\par
\editionEtape{D08 · Concours des médiatrices et cercle circonscrit}
\textbf{Indice 1.} Si deux médiatrices étaient parallèles, leurs côtés perpendiculaires le seraient aussi.\par\smallskip
\textbf{Indice 2.} Utiliser successivement $OA=OB$ et $OA=OC$.\par

\end{blocEdition}
\editionPage{indices-3}{Indications progressives}{Une aide, puis une nouvelle recherche}{Ne consulter que l’aide nécessaire.}
\begin{blocEdition}{indices-3}
\editionEtape{D09 · Un lieu parabolique}
\textbf{Indice 1.} La distance au point $F$ est $\sqrt{x^2+(y-2)^2}$.\par\smallskip
\textbf{Indice 2.} Les deux distances sont positives ou nulles : une égalité de leurs carrés est équivalente à leur égalité.\par
\editionEtape{D10 · Des inégalités décrivent une région}
\textbf{Indice 1.} Les deux droites obliques se coupent en résolvant $x+y=8$ et $x+2y=12$.\par\smallskip
\textbf{Indice 2.} Chercher des coefficients positifs $a,b$ tels que $a(x+y)+b(x+2y)=3x+4y$.\par
\editionEtape{D11 · Carré et racine, deux courbes liées}
\textbf{Indice 1.} Échanger l’abscisse et l’ordonnée ; utiliser $\sqrt{a^2}=a$ seulement parce que $a\ge0$.\par\smallskip
\textbf{Indice 2.} Comparer d’abord $x^2$ à $x$, puis appliquer le même raisonnement à la racine.\par
\editionEtape{D12 · Un paramètre et deux extrémums}
\textbf{Indice 1.} Le minimum correspond au point de $[0;4]$ le plus proche de $a$.\par\smallskip
\textbf{Indice 2.} Si les deux extrémités étaient à distance strictement inférieure à 2 de $a$, pourraient-elles rester distantes de 4 ?\par

\end{blocEdition}
\editionPage{indices-4}{Indications progressives}{Une aide, puis une nouvelle recherche}{Ne consulter que l’aide nécessaire.}
\begin{blocEdition}{indices-4}
\editionEtape{D13 · Optimiser une dispersion}
\textbf{Indice 1.} L’ordre des valeurs est connu ; la médiane est donc le terme central $t$.\par\smallskip
\textbf{Indice 2.} Utiliser $\sigma^2=\overline{x^2}-\bar x^2$ puis compléter le carré.\par
\editionEtape{D14 · Des marges qui ne fixent pas l’intersection}
\textbf{Indice 1.} Les cases sont $k$, $60-k$, $50-k$ et le reste pour arriver à 100.\par\smallskip
\textbf{Indice 2.} Chaque case doit être un entier positif ou nul.\par
\editionEtape{D15 · Défauts rares et conditionnement inversé}
\textbf{Indice 1.} Les vrais positifs viennent des $pN$ défauts ; les faux positifs des $(1-p)N$ pièces conformes.\par\smallskip
\textbf{Indice 2.} Le dénominateur est positif. Comparer la fraction à $1/2$ revient à comparer le nombre de vrais positifs à celui de faux positifs.\par
\editionEtape{D16 · Trois dés, deux sommes}
\textbf{Indice 1.} Trois nombres distincts donnent six ordres ; deux égaux, trois ordres ; trois égaux, un seul.\par\smallskip
\textbf{Indice 2.} Pour la somme 9, les triplets vont de $(1;2;6)$ à $(3;3;3)$. Pour 10, de $(1;3;6)$ à $(3;3;4)$.\par

\end{blocEdition}
